\section{Incidencia hexádico--octádica, operador de área e información sectorial}
\label{rev4:ae:seccion}

La recuperación geométrica del registro adquiere contenido físico cuando se
compone con observables cuya construcción conserva las incidencias y las
condiciones de estado utilizadas. El paso de una héxada a su óctada marcada
proporciona un caso preciso: la ampliación del soporte transforma dos
espacios de incidencia, determina sus multiplicidades y permite conservar
esa distinción en un operador de área y en el logaritmo de un estado
reducido. La exposición siguiente reúne esa construcción del segundo
artículo de la serie~\cite{hmtII} con la representación grassmanniana
marcada de la sección precedente.

El punto de partida permanece en el estado enriquecido producido por
APP--TRIT--TPK. Las dos operaciones de APP conservan sus divisiones
completas, TRIT determina régimen y orientación, y la composición
\(U_t=\operatorname{Upd}_t\circ\operatorname{Tra}_t\circ
\operatorname{Sel}_t\) transporta y actualiza la información de posición,
acarreo y memoria. La prolongación compatible y la holonomía nonádica
conservan conjuntamente las realizaciones aritmética e incidencial de la
estructura discreta del continuo. El registro \(K\), las coordenadas
regionales \(\pi_{\rm HMT},\varphi_{\rm HMT},e_{\rm HMT}\), la
coordenada \(\alpha_{\rm HMT}\) y la configuración excepcional proceden
de esa construcción común. Aquí se utilizan sus resultados y las marcas
que permiten prolongarlos hacia los observables de frontera.

\subsection{Elevación de incidencia y diferencia de multiplicidades}

Sea \(\mathcal H\) el sistema de héxadas de Witt sobre las doce posiciones
del registro. Fijemos la coordenatización de incidencia \(\ell:[12]\to\mathscr D\),
donde \(\mathscr D\) es una dodécada del código binario de Golay extendido
\(G_{24}\), de parámetros \([24,12,8]_2\). La coordenatización transporta las
héxadas de \(\mathcal H\) al diseño residual sobre \(\mathscr D\).
Este dato de realización conserva la configuración marcada determinada
en el desarrollo excepcional. Los soportes de peso ocho de \(G_{24}\)
constituyen las óctadas del sistema \(S(5,8,24)\).

\begin{proposition}[Elevación única de una héxada residual]
\label{rev4:ae:elevacion}
Los conjuntos \(O\cap\mathscr D\), con \(O\) una óctada y
\(|O\cap\mathscr D|=6\), forman un sistema \(S(5,6,12)\).
Cada héxada de ese sistema es la intersección de \(\mathscr D\) con
una única óctada. En particular, la héxada marcada \(H\subset\mathscr D\)
tiene una elevación determinada \(O_H\) y una inclusión \(H\subset O_H\).
\end{proposition}
\begin{proof}
Los pesos del código son \(0,8,12,16,24\). Si
\(j=|O\cap\mathscr D|\), la suma binaria de los vectores de soporte
\(O\) y \(\mathscr D\) tiene peso \(20-2j\); por tanto,
\(j\in\{2,4,6\}\). Toda quíntupla \(T\subset\mathscr D\)
pertenece a una única óctada. Como su intersección con \(\mathscr D\)
contiene \(T\), esa intersección tiene seis elementos. Se obtiene así
una única héxada residual sobre cada quíntupla. Si dos óctadas elevaran
la misma héxada, ambas contendrían cualquiera de sus quíntuplas y
coincidirían por la propiedad de Steiner.
\end{proof}

En lo que sigue, \(\binom H2\) denota el conjunto de subconjuntos de dos
elementos de \(H\). Conservamos ese segundo factor al ampliar el primero:
\begin{equation}
 \mathcal F_6=H\times\binom H2,
 \qquad
 \mathcal F_8=O_H\times\binom H2,
 \qquad
 \jmath(h,P)=(h,P).
 \label{rev4:ae:incidencias}
\end{equation}
Un elemento de estos espacios es una incidencia punto--pareja. El punto
distinguido recorre primero la héxada y después su óctada, mientras la
pareja sigue perteneciendo a la héxada de origen. Esta especificación
fija la operación responsable del incremento.

\begin{proposition}[Ampliación exacta del espacio de incidencias]
\label{rev4:ae:incremento}
La aplicación \(\jmath\) es inyectiva y su complemento es
\((O_H\setminus H)\times\binom H2\). En consecuencia,
\begin{equation}
 |\mathcal F_6|=6\binom62=90,
 \qquad |\mathcal F_8|=8\binom62=120,
 \qquad |\mathcal F_8\setminus\jmath(\mathcal F_6)|=30.
 \label{rev4:ae:cardinales}
\end{equation}
\end{proposition}
\begin{proof}
La inclusión \(H\subset O_H\) conserva ambas coordenadas y prueba la
inyectividad. La partición disjunta \(O_H=H\sqcup(O_H\setminus H)\)
induce la partición del producto cartesiano. Como
\(|O_H\setminus H|=2\) y \(\binom62=15\), el complemento tiene
\(2\cdot15=30\) elementos.
\end{proof}

Las dos multiplicidades se obtienen antes de cualquier distribución
probabilística sobre la frontera. La diferencia normalizada
\((90-120)/(24\binom62)=-1/12\) conserva el signo de la comparación
orientada. Su interpretación depende del espacio de incidencias en que
se ha efectuado la operación.

\subsection{Transporte colectivo y parámetro de Barbero--Immirzi}

La representación de frontera retiene las etiquetas de ambos sectores.
Para explicitar su soporte, escribamos
\(x_j=r_{2j-1}+3r_{2j}\), con \(r_k\in\{0,1,2\}\), en tres
coordenadas de \(\mathbb Z/9\mathbb Z\). El reagrupamiento posicional
\begin{equation}
 \beta_{729}(x_1,x_2,x_3)
 =\bigl(r_1+3r_2+9r_3,\ r_4+3r_5+9r_6\bigr)
 \label{rev4:ae:reagrupamiento}
\end{equation}
es una biyección hacia \((\mathbb Z/27\mathbb Z)^2\). La unicidad de
la expansión ternaria prueba que su inversa extrae los seis dígitos y
restituye las tres parejas iniciales. Se preserva el orden posicional
con su información de acarreo; las respectivas adiciones se mantienen
definidas en sus propios grupos.

En el bloque \(B_\partial=\{0,\ldots,8\}^2\) del toro de orden
veintisiete se distinguen los enlaces orientados
\begin{align}
 E_{90}&=\{((-1,j),(0,j)),\ ((8,j),(9,j)):0\le j\le8\},\\
 E_{120}&=\{((i,-1),(i,0)),\ ((i,8),(i,9)):0\le i\le8\}.
 \label{rev4:ae:enlaces}
\end{align}
Las coordenadas se interpretan módulo veintisiete. Los dos conjuntos
son disjuntos y contienen dieciocho enlaces cada uno: dos lados de
nueve enlaces por dirección. Las etiquetas \(90\) y \(120\) proceden
de los espacios de incidencia; los cardinales de los conjuntos de
enlaces son \(18\) y \(18\).

Sean \(u_m\) los vectores uniformes normalizados de
\(\ell^2(\mathcal F_6)\oplus\ell^2(\mathcal F_8)\), con
\(m=90,120\), y sean \(v_m\) los correspondientes vectores
uniformes normalizados de \(\ell^2(E_{90})\oplus\ell^2(E_{120})\).
Por ejemplo, \(u_{90}=90^{-1/2}\sum_{f\in\mathcal F_6}(\delta_f,0)\)
y \(v_{90}=18^{-1/2}\sum_{e\in E_{90}}(\delta_e,0)\).
Cada pareja es ortonormal.

\begin{proposition}[Conservación de la estructura espectral colectiva]
\label{rev4:ae:colectivo}
La aplicación lineal \(\mathcal J u_m=v_m\) es una isometría
sobreyectiva entre los dos planos colectivos. Si
\[
 D_{\rm inc}=90|u_{90}\rangle\langle u_{90}|
             +120|u_{120}\rangle\langle u_{120}|,
 \qquad
 D_\partial=90|v_{90}\rangle\langle v_{90}|
             +120|v_{120}\rangle\langle v_{120}|,
\]
entonces \(\mathcal J D_{\rm inc}=D_\partial\mathcal J\).
Los proyectores sectoriales y sus combinaciones orientadas se
transportan mediante la misma aplicación.
\end{proposition}
\begin{proof}
La aplicación lleva una base ortonormal a otra y queda determinada
por esas imágenes. La igualdad de operadores vale sobre \(u_m\),
pues ambos miembros tienen valor \(m v_m\). Los proyectores son
\((120I-D)/30\) y \((D-90I)/30\) en cada plano, de modo que
sus combinaciones se conservan por composición lineal.
\end{proof}

El transporte anterior actúa sobre los modos uniformes. Las fluctuaciones
intrínsecas de cada espacio de incidencias permanecen en sus complementos
ortogonales. La ponderación areal utiliza precisamente las dos etiquetas
que esta isometría conserva.

El parámetro de Barbero--Immirzi se determina a partir de esos datos y
de la orientación angular previamente construida. Sean \(A\) y \(C^*\)
las coordenadas angulares en grados, en el levantamiento fijado en la
sección precedente, y pongamos
\[
 x=\frac{\pi_{\rm HMT}A}{180},\qquad
 y=\frac{\pi_{\rm HMT}C^*}{180},\qquad
 q_\pm=e^{-x\mp y},\qquad
 S_m(q)=\frac{q^m}{1-q^{3m}}.
\]
Se conserva la cámara \(x>|y|\), por lo que \(0<q_\pm<1\).
La cadena orientada de una vuelta asigna doce sectores a
\(\mathcal F_6\) y una costura de retorno con orientación contraria
a \(\mathcal F_8\). Su evaluación es \(12S_{90}-S_{120}\).
El coeficiente de \((1-q)^{-1}\) queda determinado por
\[
 \lim_{q\uparrow1}(1-q)(12S_{90}(q)-S_{120}(q))
 =\frac{12}{270}-\frac1{360}=\frac1{24},
\]
porque \(1-q^{3m}=(1-q)\sum_{j=0}^{3m-1}q^j\).
La evaluación conjugada de la cadena, junto con la componente angular,
define el parámetro estructural
\begin{equation}
 \gamma_{\rm HMT}
 =\frac{\pi_{\rm HMT}AC^*}{180}
 +12\bigl[S_{90}(q_-)-S_{90}(q_+)\bigr]
 -\bigl[S_{120}(q_-)-S_{120}(q_+)\bigr].
 \label{rev4:ae:gamma}
\end{equation}
El origen de los coeficientes \(12:-1\) reside en la cadena orientada.
La factorización precedente determina su coeficiente de polo.
La inversión \(C^*\mapsto-C^*\) intercambia los canales y cambia
el signo de \(\gamma_{\rm HMT}\). Conservamos a continuación la
rama \(C^*>0\), \(\gamma_{\rm HMT}>0\), con la normalización
areal fijada por esa misma orientación.

\subsection{Retorno longitudinal y determinación de la escala areal}
\label{rev5:longitud}

La unidad dimensional del área admite una construcción explícita
desde el mismo estado enriquecido. Escribiremos \(L_\alpha\) para la
longitud física obtenida; el entero \(L\) que designa la longitud de un
registro en la sección~\ref{sec:compat-lectores-formas} conserva su
significado combinatorio. La distinción separa el cardinal de una
representación, la extensión de una arista y el radio de su realización.

La composición utiliza las coordenadas regionales y el registro ya
generados por APP--TRIT--TPK. APP conserva ambas hojas con sus residuos
y cocientes; TRIT orienta el régimen local; TPK prolonga el estado,
conservando acarreo, frontera, ruta y memoria. La cadena
\(w_6\to w_{12}\to w_{18}\to w_{24}\to w_{30}\to R_{36}\to G_9\)
y el orden \(U_t\to\Gamma_9\to K_{\rm ph}\) mantienen el
retorno de fase con incremento de memoria. La construcción conjunta
del continuo conserva sus aspectos espectral, solenoidal, gauge,
cohomológico y determinantal. La longitud que sigue es una lectura
dimensional posterior de esa estructura y de sus salidas correlacionadas
\(\pi_{\rm HMT},\varphi_{\rm HMT},e_{\rm HMT},\alpha_{\rm HMT}\).

Fijemos el registro marcado
\[
 \Omega=\{(j,k,q,u,v):j,k\in\mathbb Z/12\mathbb Z,\quad
 k-j\in\{0,3,6,9\},\quad1\le q\le8,\quad
 u<v,\quad u,v\in\{1,2,4,5,7,8\}\}.
\]
La resta se calcula en ese grupo cíclico. Su cardinal es
\(N_\Omega=12\cdot4\cdot8\binom62=5760\).
El factor \(120\) procede del espacio octádico de incidencias
punto--pareja ya construido; las doce posiciones y las cuatro
posiciones de cada órbita conservan sus marcas respectivas. Esta
resolución amplía el portador del retorno de Hadamard. En los espacios
hermitianos de partida y llegada, de dimensión \(N_\Omega\), sea
\(B\) el inversor transportado, con
\(B^*B=BB^*=I/4\). La identidad local
\(H_4^*H_4=4I\) produce esa normalización para el inversor;
su extensión ortogonal y el transporte de las marcas la conservan.

En el espacio de llegada, sean
\(u_\Omega=N_\Omega^{-1/2}\mathbf1\),
\(P=u_\Omega u_\Omega^*\), \(C=(I-P)B\) y \(M=PB\).
Se conserva el transporte completo \(C+M=B\). Los proyectores
atómicos \(P_i=e_ie_i^*\), determinados por la inscripción,
definen la expectativa diagonal
\(\Delta_\Omega(T)=\sum_iP_iTP_i\).

\begin{proposition}[Respuesta registrada y complemento conservado]
\label{rev5:retorno}
Entre las aplicaciones lineales hacia el álgebra diagonal que fijan
los \(P_i\) y son bimodulares respecto de esa álgebra,
\(\Delta_\Omega\) es única. Además,
\begin{equation}
 \Delta_\Omega(CC^*)=r_\Omega I,
 \qquad r_\Omega=\frac{N_\Omega-1}{4N_\Omega}
 =\frac{5759}{23040},
 \qquad \Delta_\Omega(MM^*)=\frac{I}{4N_\Omega}.
 \label{rev5:retorno-escalar}
\end{equation}
Las dos respuestas suman \(I/4\).
\end{proposition}
\begin{proof}
Para la unidad matricial \(E_{ij}\), la bimodularidad impone
\(E(E_{ij})=P_iE(E_{ij})P_j\). Como la imagen es diagonal,
este término vale cero si \(i\ne j\); si \(i=j\), fijar
\(P_i\) determina su imagen. Las unidades matriciales constituyen
una base y determinan el mapa. Por otra parte,
\(CC^*=(I-P)/4\), \(MM^*=P/4\), y cada entrada diagonal
de \(P\) es \(1/N_\Omega\). Aplicar la expectativa demuestra
las tres identidades.
\end{proof}

La inscripción isométrica \(We_i=e_i\otimes e_i\) realiza
\(W^*(T\otimes I)W=\Delta_\Omega(T)\). La inscripción
completa conserva la información del registro; la expectativa extrae
su respuesta diagonal. El complemento \(M\) permanece en el
transporte y posee la respuesta específica de
\eqref{rev5:retorno-escalar}. Esta separación identifica el observable
que interviene en la longitud y el componente complementario conservado.

La forma normal de Smith de \(H_4\) es
\(\operatorname{diag}(1,2,2,4)\); su conúcleo tiene orden
dieciséis. El producto del escalar generado
\(\alpha=\alpha_{\rm HMT}\), indexado por las clases de ese
conúcleo, define el lector multiplicativo finito
\[
 n_H:=\prod_{\bar v\in\operatorname{coker}_{\mathbb Z}(H_4)}
 \alpha=\alpha^{16}.
\]
Sea \(L_*>0\) una sección de referencia
radial de la recta dimensional de longitud. La composición longitudinal
actúa linealmente sobre una cocadena orientada \(\beta_*\):
\begin{equation}
 \mathscr T_L\beta_*
 =\alpha^{16}(\Delta_\Omega(CC^*)\otimes I)\beta_*
 =q_\alpha\beta_*,\qquad
 q_\alpha=\alpha^{16}r_\Omega,\qquad
 L_\alpha=q_\alpha L_*.
 \label{rev5:longitud-generada}
\end{equation}
La fórmula fija una composición concreta de lectores. La linealidad
conserva el tipo de longitud y su orientación; el coeficiente procede
del retorno registrado y de la norma local. Por ejemplo, la respuesta
diagonal de una fuente \(a e_i\) es \(a r_\Omega\), mientras
la norma \(\|C^*ae_i\|=|a|\sqrt{r_\Omega}\) constituye
otra lectura del mismo transporte.

Si \(\beta_*(\bar a)=-U(a)\beta_*(a)\), multiplicar por
\(q_\alpha\) conserva la inversión orientada. Si una arista se
subdivide en nueve tramos compatibles,
\(\beta_*(a)=\sum_jU_j^{-1}\beta_*(a_j)\), la misma
multiplicación conserva esa identidad y la concatenación de memoria.
En una ampliación auxiliar del registro, se transportan
\(P\), \(CC^*\) y \(\Delta_\Omega\) por tensorización
con la identidad. Así permanece \(r_\Omega\); el cardinal de
una malla refinada y el cardinal del registro marcado tienen funciones
distintas.

\subsection{Reciprocidad radial y reconocimiento de la longitud de Planck}
\label{rev5:radial}

La lectura circular del calendario de ciento ocho avances conserva
\(108\ell_0=2\pi_{\rm HMT}L_\alpha\), con
\(\ell_0=ct_0\). Por consiguiente,
\begin{equation}
 t_0=\frac{\pi_{\rm HMT}L_\alpha}{54c},\qquad
 \omega=\frac{c}{L_\alpha},\qquad
 Q_{\rm clk}=\hbar_{\rm ret}\omega
 =\frac{\hbar_{\rm ret}c}{L_\alpha}.
 \label{rev5:reloj}
\end{equation}
La acción retornada se recibe de su ecuación HMT completa en
\eqref{eq:mass-k-action}; la velocidad se recibe de los dos canales
constitutivos de la sección~\ref{sec:compat-lectores-formas}.
La relación circular convierte el avance levantado en radio. Si
\(c=\widehat c\,u_C\), \(t_0=u_T\) y
\(u_L=u_Cu_T\), resulta
\(L_*/u_L=54\widehat c/(\pi_{\rm HMT}q_\alpha)\).
La referencia radial \(L_*\) y la base de arista \(u_L\)
quedan relacionadas con sus tipos dimensionales explícitos.

En una fibra finita, sea \(X=X^*>0\) un carácter positivo
invertible, \(U=2B\) y \(Y=UXU^*\). El transporte
longitudinal es \(D_L=L_\alpha U\). Su respuesta radial positiva
se define mediante la métrica de llegada,
\(\|R_Xv\|^2=\|XD_L^*v\|^2\) para todo \(v\).
La energía y la longitud conjugada conservan el mismo carácter:
\[
 H_X=Q_{\rm clk}Y,\qquad M_X=H_X/c^2,\qquad
 \Lambda_X=L_\alpha Y^{-1}.
\]

\begin{theorem}[Coeficiente radial común y escala areal]
\label{rev5:rigidez-radial}
Las condiciones anteriores determinan una única respuesta positiva
\(R_X=L_\alpha Y\). Se cumplen
\begin{equation}
 R_X\Lambda_X=L_\alpha^2I,\qquad
 H_X\Lambda_X=\hbar_{\rm ret}cI,\qquad
 R_X=\frac{L_\alpha^2}{\hbar_{\rm ret}c}H_X.
\end{equation}
En la realización física que identifica \(R_X\) con el radio
gravitatorio reducido \(R_X=(G/c^2)M_X\), el coeficiente es
único y común a todos los caracteres positivos admisibles:
\begin{equation}
 G=\frac{c^3L_\alpha^2}{\hbar_{\rm ret}},\qquad
 \frac{\hbar_{\rm ret}G}{c^3}=L_\alpha^2.
 \label{rev5:G-area}
\end{equation}
\end{theorem}
\begin{proof}
La polarización de la igualdad de normas determina
\(R_X^2=D_LX^2D_L^*=L_\alpha^2UX^2U^*\).
El operador \(L_\alpha UXU^*\) es positivo y su cuadrado es
el indicado. La unicidad de la raíz positiva prueba la primera
afirmación. Las dos identidades de productos siguen por sustitución.
La convención física da
\(L_\alpha Y=(GQ_{\rm clk}/c^4)Y\). La invertibilidad
de \(Y\) permite cancelarlo, y la expresión de
\(Q_{\rm clk}\) proporciona \eqref{rev5:G-area}.
Si otro coeficiente conserva los mismos datos, su diferencia
multiplicada por \(M_X\) vale cero y, por tanto, es nula.
\end{proof}

El dominio de la prueba es la fibra positiva finita; también admite
operadores acotados uniformemente positivos con las mismas identidades.
La lectura de radio gravitatorio reducido constituye la convención
física declarada de esta realización. El reconocimiento convencional
posterior \(\ell_P^2=\hbar_{\rm ret}G/c^3\) identifica
\(\ell_P=L_\alpha\). La longitud se obtiene primero mediante
\eqref{rev5:longitud-generada}; esta identificación conserva su
constructor. La fórmula circular equivalente es
\[
 G=\left(\frac{54}{\pi_{\rm HMT}}\right)^2
 \frac{c^5t_0^2}{\hbar_{\rm ret}}.
\]
El factor circular ya está incluido cuando se utiliza \(L_\alpha\).
La longitud y la sección de acción determinan asimismo
\[
 t_P=\frac{L_\alpha}{c},\qquad
 m_P=\frac{\hbar_{\rm ret}}{cL_\alpha},\qquad
 E_P=\frac{\hbar_{\rm ret}c}{L_\alpha}.
\]
Estas identidades siguen al sustituir \eqref{rev5:G-area} en las
expresiones convencionales de las tres escalas; la positividad
determina sus raíces. Para un modo de carácter \(y>0\), las
lecturas son \(m(y)=m_Py\), \(r_g(y)=L_\alpha y\) y
\(\bar\lambda(y)=L_\alpha/y\). Su producto longitudinal
vale \(r_g(y)\bar\lambda(y)=L_\alpha^2\). La igualdad de
ambas longitudes equivale a \(y=y^{-1}\) y, por positividad,
caracteriza exactamente \(y=1\). El calendario conserva
\(t_0=(\pi_{\rm HMT}/54)t_P\) y
\(108t_0=2\pi_{\rm HMT}t_P\).

Un radio de horizonte \(R_h=L_\alpha\zeta_h\) conserva su
factor propio \(\zeta_h=R_h/L_\alpha\). Su área espacial
pertenece al horizonte declarado; la escala elemental del operador
sectorial siguiente es \(L_\alpha^2\). En la especialización
de Schwarzschild para la misma masa, el radio es \(2r_g(y)\)
y, por tanto, \(\zeta_h=2y\); un horizonte de otra realización
conserva su condición geométrica de formación.


\subsection{Ocupaciones de frontera y reconstrucción sectorial}

Cada enlace de \(E_{90}\sqcup E_{120}\) determina un factor
\(\mathbb C^3\), con base \(|0\rangle,|1\rangle,|2\rangle\)
y operador número \(N_e=\operatorname{diag}(0,1,2)\). El espacio de
estados de frontera es
\(\mathscr H_A=(\mathbb C^3)^{\otimes18}\otimes
(\mathbb C^3)^{\otimes18}\). Definimos \(N_m\) como la suma de
los operadores número de los dieciocho enlaces del sector \(m\).
Los dos operadores conmutan y cada uno tiene espectro entero entre
cero y treinta y seis.

La escala angular procede del reloj y del carácter normalizado del
triplete conjugado:
\begin{equation}
 \theta_{\rm clk}=\frac{C^*}{6}\frac{\pi_{\rm HMT}}{180},
 \quad
 \chi_{10}=\frac13\Tr\operatorname{diag}
 (1,e^{i\pi_{\rm HMT}/18},e^{-i\pi_{\rm HMT}/18})
 =\frac{1+2\cos(\pi_{\rm HMT}/18)}3,
 \quad a_0=\chi_{10}\theta_{\rm clk}.
 \label{rev4:ae:escala}
\end{equation}
La identidad del carácter se obtiene sumando las fases conjugadas.
En la rama considerada \(a_0>0\). La unidad areal es
\(L_\alpha^2=\ell_P^2\), construida y reconocida en
\eqref{rev5:longitud-generada}--\eqref{rev5:G-area}. Escribamos
\(a_\partial=\gamma_{\rm HMT}a_0>0\). El operador de área es
\begin{equation}
 \widehat{\mathcal A}=L_\alpha^2a_\partial N,
 \qquad N=N_{90}+N_{120},
 \qquad D=90N_{90}+120N_{120}.
 \label{rev4:ae:area}
\end{equation}
Su espectro es \(\{nL_\alpha^2a_\partial:0\le n\le72\}\).
La multiplicidad de \(n\) es el coeficiente de \(z^n\) en
\((1+z+z^2)^{36}\), pues cada factor tensorial aporta una de
las tres ocupaciones.

\begin{theorem}[Reconstrucción de las dos ocupaciones sectoriales]
\label{rev4:ae:recuperacion}
La pareja de observables \((N,D)\) determina
\begin{equation}
 N_{90}=4N-\frac{D}{30},
 \qquad N_{120}=\frac{D}{30}-3N.
 \label{rev4:ae:inversa}
\end{equation}
En el espectro conjunto, sus valores \((n,d)\) satisfacen
\(n\in\mathbb Z\), \(d\in30\mathbb Z\), \(90n\le d\le120n\) y
\(0\le4n-d/30\le36\), \(0\le d/30-3n\le36\).
Estas condiciones caracterizan la imagen de las ocupaciones enteras.
\end{theorem}
\begin{proof}
La transformación de \((N_{90},N_{120})\) en \((N,D)\) tiene matriz
\(\left(\begin{smallmatrix}1&1\\90&120\end{smallmatrix}\right)\),
cuyo determinante es \(30\). Su inversa da las identidades de
operadores. La conmutatividad permite aplicarlas en cada autoespacio
conjunto. Las desigualdades expresan la no negatividad y las capacidades
máximas de los sectores; la divisibilidad hace enteras ambas ocupaciones.
Recíprocamente, cada entero entre cero y treinta y seis se expresa como
suma de dieciocho términos de \(\{0,1,2\}\), y por tanto todo par
que satisface esas condiciones pertenece al espectro conjunto.
\end{proof}

El área conserva la suma de ocupaciones; el observable ponderado conserva
su composición sectorial. Así, \((N_{90},N_{120})=(1,0)\) y \((0,1)\)
tienen igual área \(L_\alpha^2a_\partial\) y pesos \(90\) y \(120\).
En sentido complementario, \((4,0)\) y \((0,3)\) tienen igual peso
\(D=360\), pero áreas diferentes. La pareja separa ambas situaciones.
Dentro de cada autoespacio conjunto se conservan las degeneraciones
microscópicas correspondientes a las distribuciones entre enlaces.

\begin{figure}[htbp]
\centering
\begin{tikzpicture}[>=Latex,every node/.style={font=\small},node distance=15mm]
\node[draw,rounded corners,align=center] (a) {Ocupación sectorial\\$(1,0)$};
\node[draw,rounded corners,align=center,right=54mm of a] (b) {Ocupación sectorial\\$(0,1)$};
\node[draw,rounded corners,align=center] (n) at ($(a)!0.5!(b)+(0,-1.8)$)
 {Área común\\$L_\alpha^2a_\partial$};
\draw[->] (a)--node[left] {$N=1$}(n);
\draw[->] (b)--node[right] {$N=1$}(n);
\node[above=5mm of a] {$D=90$};
\node[above=5mm of b] {$D=120$};
\end{tikzpicture}
\caption{La igualdad de área conserva la ocupación total. La ponderación
incidencial distingue las dos composiciones; su lectura conjunta permite
reconstruir las ocupaciones de ambos sectores.}
\label{rev4:ae:figura}
\end{figure}

\subsection{Estado reducido, purificación y entropía de entrelazamiento}

La interpretación informacional requiere fijar el estado. Para
\(\Delta\in\mathbb R\) finito, con la normalización del parámetro
modular de la construcción sectorial, definimos
\begin{equation}
 x_m=e^{-m\Delta},\qquad z_m=1+x_m+x_m^2,
 \qquad Z=z_{90}^{18}z_{120}^{18},
 \qquad \rho_A(\Delta)=Z^{-1}e^{-\Delta D}.
 \label{rev4:ae:rho}
\end{equation}
La suma de ocupaciones conmutantes y la descomposición tensorial prueban
la expresión de \(Z\). Todos los autovalores de \(\rho_A\) son
estrictamente positivos para \(\Delta\) real finito. Su Hamiltoniano
modular, denotado \(K_A\) para distinguirlo del registro aritmético
\(K\), es
\begin{equation}
 K_A=-\log\rho_A=(\log Z)I+\Delta D.
 \label{rev4:ae:modular}
\end{equation}
Este operador caracteriza la distribución espectral del estado reducido.
El generador de evolución temporal se define en el desarrollo dinámico
con sus propios parámetros y dominio. Para \(\Delta\ne0\), conocido
\(\Delta\) y conocida la normalización, la pareja
\((\widehat{\mathcal A},K_A)\) recupera \((N,D)\) y, mediante
\eqref{rev4:ae:inversa}, ambas ocupaciones. Para \(\Delta=0\),
el estado es uniforme y \(K_A=36\log3\,I\).

Construimos la bipartición de manera explícita. Sea \(\mathscr H_{\bar A}\)
una copia de \(\mathscr H_A\), con las bases conjugadas fijadas. Para
un enlace de peso \(w\), el vector
\[
 |\operatorname{TFD}(w)\rangle
 =\frac1{\sqrt{1+e^{-w}+e^{-2w}}}
 \sum_{n=0}^{2}e^{-wn/2}|n\rangle_A\otimes|n\rangle_{\bar A}
\]
tiene norma uno. El estado puro
\begin{equation}
 |\Psi_\Delta\rangle
 =|\operatorname{TFD}(90\Delta)\rangle^{\otimes18}
  \otimes|\operatorname{TFD}(120\Delta)\rangle^{\otimes18}
 \label{rev4:ae:tfd}
\end{equation}
determina la bipartición \(A\mid\bar A\).

\begin{proposition}[Reducción sectorial y entropía de la purificación]
\label{rev4:ae:entropia}
La traza parcial de \(|\Psi_\Delta\rangle\langle\Psi_\Delta|\)
sobre \(\mathscr H_{\bar A}\) es \(\rho_A(\Delta)\). Definidos
\[
 f_m=\frac{x_m+2x_m^2}{z_m},
 \qquad s(\Delta)=-\Tr(\rho_A(\Delta)\log\rho_A(\Delta)),
\]
se cumplen
\begin{align}
 \langle N\rangle_\Delta&=18(f_{90}+f_{120}),\\
 \langle D\rangle_\Delta&=18(90f_{90}+120f_{120}),\\
 \langle\widehat{\mathcal A}\rangle_\Delta
 &=18L_\alpha^2a_\partial(f_{90}+f_{120}),\\
 s(\Delta)&=18\log z_{90}+18\log z_{120}
             +\Delta\langle D\rangle_\Delta.
 \label{rev4:ae:entropia-formula}
\end{align}
La última expresión es la entropía de entrelazamiento de la
purificación \eqref{rev4:ae:tfd}, en unidades adimensionales.
\end{proposition}
\begin{proof}
En la traza parcial de un enlace, la ortogonalidad de
\(|n\rangle_{\bar A}\) elimina los términos de índices distintos y
deja los pesos \((1,x_m,x_m^2)/z_m\). La traza parcial de un producto
tensorial es el producto de las trazas parciales; se obtiene
\eqref{rev4:ae:rho}. El primer momento de cada enlace es \(f_m\).
La aditividad de los operadores número da las tres primeras identidades.
La última resulta de sustituir \eqref{rev4:ae:modular} en
\(s=\Tr(\rho_AK_A)\). Los coeficientes de la construcción anterior
son coeficientes de Schmidt de la bipartición declarada, por lo que
esa entropía coincide con su entropía de entrelazamiento.
\end{proof}

La entropía máxima de este espacio es \(36\log3\), alcanzada en
\(\Delta=0\). La cantidad \(36\log3-s(\Delta)\) es la
entropía relativa de \(\rho_A(\Delta)\) respecto de
\(I/3^{36}\). La incidencia fija los sectores y sus multiplicidades;
el parámetro del estado fija las probabilidades; la coordenatización dimensional
convierte la ocupación en área. La identificación física de
\(\mathscr H_{\bar A}\) conserva la bipartición y el álgebra de
observables que acaban de especificarse.

\subsection{Identidad modular finita y transporte de la información geométrica}

Fijemos el estado de referencia \(\rho_0=\rho_A(\Delta_0)\), con
\(\Delta_0\) real finito, y mantengamos constantes
\(\Delta_0,a_\partial,L_\alpha\). Las áreas sectoriales
normalizadas \(\mathcal A_m^{\rm n}=a_\partial N_m\) satisfacen
\begin{equation}
 K_0=-\log\rho_0=c_0I+
 \beta_{90}\mathcal A_{90}^{\rm n}
 +\beta_{120}\mathcal A_{120}^{\rm n},
 \qquad c_0=\log Z(\Delta_0),\qquad
 \beta_m=\frac{m\Delta_0}{a_\partial}.
 \label{rev4:ae:betas}
\end{equation}
La relación \(\beta_{120}=\tfrac43\beta_{90}\) conserva la razón
de incidencias. Para \(\Delta_0\ne0\), la razón de ambos coeficientes
está definida y vale \(4/3\).

\begin{theorem}[Identidad finita de área e información relativa]
\label{rev4:ae:ley-finita}
Para todo estado normalizado \(\sigma\) de \(\mathscr H_A\), sea
\(\delta_\sigma\langle B\rangle=\Tr[(\sigma-\rho_0)B]\).
Entonces
\begin{equation}
 s(\sigma)-s(\rho_0)
 =\sum_{m\in\{90,120\}}\beta_m
       \delta_\sigma\langle\mathcal A_m^{\rm n}\rangle
   -D(\sigma\Vert\rho_0),
 \label{rev4:ae:ley-finita-formula}
\end{equation}
donde \(D(\sigma\Vert\rho_0)=
\Tr[\sigma(\log\sigma-\log\rho_0)]\) es finita y no negativa,
con la convención \(0\log0=0\).
\end{theorem}
\begin{proof}
De la definición se obtiene
\(D(\sigma\Vert\rho_0)=-s(\sigma)+\Tr(\sigma K_0)\),
mientras \(s(\rho_0)=\Tr(\rho_0K_0)\). Restar las identidades
y sustituir \eqref{rev4:ae:betas} prueba
\eqref{rev4:ae:ley-finita-formula}; el término escalar desaparece
por la normalización de las trazas. Para la positividad, sean \(p_i\)
y \(u_i\) los autovalores y autovectores de \(\sigma\), y
\(q_j>0,v_j\) los de \(\rho_0\). Con
\(r_i=\langle u_i,\rho_0u_i\rangle\), la concavidad de
\(\log\) da
\(\langle u_i,\log\rho_0u_i\rangle\le\log r_i\).
Así, \(D(\sigma\Vert\rho_0)\ge\sum_i p_i\log(p_i/r_i)\ge0\),
por convexidad de \(t\log t\), ya que \(\sum_i r_i=1\).
El rango completo de \(\rho_0\) garantiza la finitud.
\end{proof}

Su versión infinitesimal mantiene fijo el estado de referencia. Si
\(\sigma(\varepsilon)\) es diferenciable, normalizado y
\(\sigma(0)=\rho_0\), entonces
\[
 \left.\frac{\dd s(\sigma(\varepsilon))}{\dd\varepsilon}\right|_0
 =\sum_m\beta_m
  \left.\frac{\dd\langle\mathcal A_m^{\rm n}\rangle_{
                          \sigma(\varepsilon)}}{\dd\varepsilon}\right|_0.
\]
En efecto, la derivada de la traza de \(-\sigma\log\sigma\) es
\(-\Tr[\dot\sigma(\log\rho_0+I)]\); la normalización elimina
\(\Tr\dot\sigma\) y la sustitución de \(K_0\) da la fórmula.
La identidad finita añade a esa respuesta tangente la corrección exacta
de entropía relativa. Los coeficientes de las áreas físicas son
\(\beta_m/L_\alpha^2\), conservando la misma coordenatización dimensional.

La composición con la representación grassmanniana precisa finalmente
qué información se conserva. En
\(\mathcal F(K,\xi)=([B(K)],Q,\xi)\), las coordenadas adicionales
\(\xi\) retienen las magnitudes regionales, la orientación angular,
la inclusión marcada, la identificación de los sectores, la base de
ocupaciones, el parámetro \(\Delta\) y la unidad areal pertinente.
La inversa de la proposición~\ref{prop:compat-inversa-marcada}
recupera \(K\); las otras condiciones se conservan mediante sus
marcas. Por el teorema~\ref{thm:compat-algebra-lectores}, la composición
con esa inversa transporta \(\gamma_{\rm HMT}\),
\(\widehat{\mathcal A}\), \(\rho_A\), \(K_A\) y sus evaluaciones
en el dominio declarado. Los cambios de coordenatización del mismo punto marcado
conservan esos objetos, y los refinamientos conservan las evaluaciones
heredadas cuando mantienen los mismos datos sectoriales.

Esta articulación conserva tres niveles distinguibles de información.
La representación positiva marcada recupera el registro con los datos
retenidos; la pareja área--Hamiltoniano modular recupera las dos
ocupaciones bajo las condiciones de \eqref{rev4:ae:modular}; la
purificación determina el entrelazamiento de una bipartición concreta.
La pérdida de información en una proyección se calcula por sus fibras:
el área identifica ocupaciones de igual suma, mientras la lectura
conjunta separa su composición y conserva la degeneración microscópica
restante. El vínculo entre geometría e información queda así expresado
por aplicaciones e identidades verificables sobre la misma construcción.
\subsection{Naturalidad de la escala y formas canónicas operatoriales}
\label{rev5:naturalidad-areal}

La longitud generada permite precisar la compatibilidad entre la
realización positiva, el operador de área y la información sectorial.
La representación marcada \(\mathcal F(K,\xi)\) conserva,
entre los datos pertinentes de \(\xi\), la resolución
\(\Omega\), sus proyectores, el retorno transportado y la sección
\(L_*\). Las coordenadas regionales generadas, la acción y los canales
constitutivos mantienen su genealogía anterior. La inversa marcada
restituye así todos los argumentos de
\eqref{rev5:longitud-generada} y \eqref{rev5:G-area}.
En particular, el lector geométrico ya construido de
\(\alpha^{\mathrm{geom}}\) determina las expresiones
\[
 L_\alpha^{\mathrm{geom}}
 =L_*(\alpha^{\mathrm{geom}})^{16}\frac{5759}{23040},
 \qquad
 G^{\mathrm{geom}}
 =\frac{(c^{\mathrm{geom}})^3
 (L_\alpha^{\mathrm{geom}})^2}{\hbar_{\rm ret}^{\mathrm{geom}}}.
\]
Los lectores de velocidad y acción se obtienen por la misma
composición con la inversa marcada. Estas expresiones conservan
los datos del retorno y la sección radial de referencia.

\begin{corollary}[Conservación por cambios de coordenatización y refinamiento heredado]
\label{rev5:conservacion}
Los cambios positivos de coordenatización que conservan el punto marcado y los
refinamientos que restituyen el registro inicial conservan
\(L_\alpha\), \(G\), \(\gamma_{\rm HMT}\) y
\(\widehat{\mathcal A}=L_\alpha^2a_\partial N\), cuando
transportan asimismo los datos sectoriales y dimensionales declarados.
El espectro normalizado
\(\widehat{\mathcal A}/L_\alpha^2=a_\partial N\), sus
multiplicidades y el estado \(\rho_A(\Delta)\) permanecen iguales.
\end{corollary}
\begin{proof}
El teorema~\ref{thm:compat-algebra-lectores} transporta cada lector
por composición con la inversa. Las mutaciones del mismo punto
reconstruyen el mismo registro. En el refinamiento heredado, sumar
las separaciones subdivididas reconstruye cada coordenada inicial.
La tensorización de la inscripción conserva \(r_\Omega\),
y sus restantes argumentos quedan fijados por las marcas. Por
sustitución se conservan la longitud, el coeficiente gravitatorio y
el operador de área. Los proyectores de ocupación conservan el
polinomio de multiplicidades \((1+z+z^2)^{36}\) y la expresión
normalizada del estado.
\end{proof}

El refinamiento heredado considerado aquí conserva el mismo espacio
de treinta y seis enlaces de frontera. Una ampliación que introduzca
un espacio físico de frontera diferente requiere entrelazadores
especificados para \(N_{90}\) y \(N_{120}\), junto con el
transporte de los datos sectoriales declarados.

Sea ahora \(P\) un politopo de rango uno de este manuscrito,
con una subdivisión orientada \(P=\bigcup_\sigma P_\sigma\)
del mismo soporte, multiplicidad uno y caras compatibles. Se fija
la misma fibra sectorial \(\mathscr H_A\) sobre todas sus
celdas. La forma de grado máximo con valores operatoriales y
dimensión de área se define por
\[
 \boldsymbol\Omega^{\mathcal A}_P
 =\widehat{\mathcal A}\otimes\Omega_P.
\]
Aquí \(\widehat{\mathcal A}\) es constante respecto de las
coordenadas diferenciales del politopo: procede del estado HMT
fijado, mientras \(\Omega_P\) describe su soporte geométrico.
El coeficiente \(\widehat{\mathcal A}\) es semidefinido positivo.
Su componente en el sector de vacío \(N=0\) es nula. Sobre un
autoespacio de ocupación \(n>0\), dividir por
\(nL_\alpha^2a_\partial\) recupera la forma canónica
\(\Omega_P\) multiplicada por la identidad de ese autoespacio.

\begin{proposition}[Aditividad areal y residuos de frontera]
\label{rev5:forma-areal}
En la trivialización sectorial común se cumplen
\begin{equation}
 \boldsymbol\Omega^{\mathcal A}_P
 =\sum_\sigma\widehat{\mathcal A}\otimes\Omega_{P_\sigma},
 \qquad
 \operatorname{Res}_F\boldsymbol\Omega^{\mathcal A}_P
 =\widehat{\mathcal A}\otimes\Omega_F.
\end{equation}
Las identidades persisten bajo subdivisiones sucesivas y residuos
iterados. Para un estado fijo \(\rho\), tomar la traza da
\(\Tr(\rho\widehat{\mathcal A})\Omega_P\) y conserva
las mismas identidades escalares.
\end{proposition}
\begin{proof}
El espacio de operadores de \(\mathscr H_A\) es finito
dimensional. En cualquier base, la afirmación es la identidad de
aditividad de las formas canónicas multiplicada por cada entrada
constante de \(\widehat{\mathcal A}\). El residuo y la traza
son lineales. La cancelación de facetas interiores conserva sus
orientaciones y la misma entrada operatorial en ambos lados.
La iteración aplica esas operaciones a cada subdivisión y a cada
bandera de caras.
\end{proof}

Si distintas celdas llevan coeficientes operatoriales regulares
\(A_\sigma\), el defecto de una pareja sobre su cara común es
\((A_\sigma|_F-A_\tau|_F)\otimes\Omega_F\).
La igualdad de restricciones operatoriales a la frontera proporciona
el pegado, bajo las condiciones de polos simples del
teorema~\ref{thm:compat-pegado}. La suma global conserva además
todas las contribuciones sobre cada divisor ambiental y el control
de polos exteriores. La normalización dimensional, la cancelación
de residuos y la cobertura del soporte tienen así criterios
separados y compatibles.

La identidad modular finita también conserva una expresión directa
en áreas físicas \(\mathcal A_m=L_\alpha^2a_\partial N_m\):
\begin{equation}
 s(\sigma)-s(\rho_0)
 =\sum_{m=90,120}\frac{m\Delta_0}{L_\alpha^2a_\partial}
   \Tr[(\sigma-\rho_0)\mathcal A_m]
  -D(\sigma\Vert\rho_0).
 \label{rev5:modular-fisica}
\end{equation}
Es \eqref{rev4:ae:ley-finita-formula} después de insertar la
longitud construida. El factor angular \(a_0\), contenido en
\(a_\partial=\gamma_{\rm HMT}a_0\), permanece íntegro.
Si se comparan secciones dimensionales con \(L'_\alpha=sL_\alpha\)
y se conserva el mismo estado y los mismos datos adimensionales,
las áreas se multiplican por \(s^2\) y sus coeficientes modulares
por \(s^{-2}\). La entropía y la entropía relativa permanecen
iguales. Esta comparación de secciones no constituye otra solución
para los mismos datos longitudinales fijados.

Para un horizonte esférico de radio \(R_h=L_\alpha\zeta_h\),
el área geométrica vale \(4\pi_{\rm HMT}L_\alpha^2\zeta_h^2\).
Su identificación con la esperanza del operador sectorial, si se
requiere para un mismo corte, exige la condición explícita
\[
 a_\partial\Tr(\rho N)=4\pi_{\rm HMT}\zeta_h^2.
\]
La geometría del horizonte y el estado de frontera deben satisfacer
esa relación con sus propias construcciones. El factor
\(\zeta_h\) caracteriza ese horizonte; la escala longitudinal
del acoplamiento permanece \(L_\alpha\).

